% !TEX root = ../thesis.tex

\chapter{User Manual}\label{app:user.manual}

\begin{sloppypar}

This user manual describes how to use the implemented information output of the thesis: a reproducible Python pipeline for binary semantic segmentation of archaeological features in LiDAR-derived raster data.

The pipeline is intended for experiments with prepared raster datasets, training segmentation models, evaluating trained checkpoints, and generating full-area \gls{probabilitymap}. It is not an interactive desktop application. The expected user is therefore a researcher or student who can run Python scripts from a command line and who has access to the prepared data files.

\section{Repository and Directory Structure}

The source repository is available as:

\begin{verbatim}
https://github.com/kitnew/lidar-archaeology-segmentation.git
\end{verbatim}

The most important directories are:

\begin{description}
    \item[\texttt{thesis/src/}] Python source code for datasets, models, training, evaluation, inference, losses, and metrics.
    \item[\texttt{thesis/configs/}] Hydra configuration files defining datasets, models, losses, paths, and experiment parameters.
    \item[\texttt{thesis/data/processed/}] Expected location of prepared \texttt{.npz} raster datasets and tile \gls{datasetmanifest}. The data files are not distributed through Git because of size and privacy constraints.
    \item[\texttt{thesis/outputs/}] Default location for trained checkpoints, Hydra logs, and generated prediction files.
    \item[\texttt{thesis/notebooks/} and \texttt{thesis/notebooks/scripts/}] Helper notebooks and scripts used during data preparation and visualization.
\end{description}

\section{System Requirements}

The pipeline was developed for Linux and Python. GPU execution with CUDA is expected for training and evaluation. The training script explicitly selects \texttt{cuda}, therefore a CUDA-capable environment is required for normal training runs.

The main software requirements are Python 3.8 or newer, PyTorch, Torchvision, Hydra, OmegaConf, NumPy, Pandas, Pillow, Matplotlib, \gls{gdal}, tqdm, and segmentation-models-pytorch for DeepLabV3+ experiments. The repository contains \texttt{requirements.txt} and \texttt{pyproject.toml}; however, the exact PyTorch and CUDA builds should be installed according to the target machine.

\section{Installation}

After cloning the repository, create and activate a virtual environment:

\begin{verbatim}
python -m venv .venv
source .venv/bin/activate
pip install -r requirements.txt
\end{verbatim}

If DeepLabV3+ configurations are used, install \texttt{segmentation-models-pytorch}. If the local \texttt{requirements.txt} does not include PyTorch, install the correct PyTorch, Torchvision, and CUDA build for the machine before running the pipeline.

\section{Input Data}

The training and inference scripts expect prepared NumPy archives and tile \gls{datasetmanifest}. A prepared dataset archive normally contains:

\begin{itemize}
    \item \texttt{data}: raster input tensor with shape \texttt{(C, H, W)};
    \item \texttt{valid}: valid-data mask with shape \texttt{(H, W)};
    \item \texttt{mask}: binary reference mask with shape \texttt{(H, W)};
    \item \texttt{soft\_mask}: optional soft reference mask used for training when enabled;
    \item \texttt{meta}: optional geospatial metadata.
\end{itemize}

The \gls{datasetmanifest} file, usually \texttt{.parquet}, contains the tile coordinates, tile size, subset assignment, validity fraction, and whether the tile contains positive pixels. Dataset variants are selected through files in \texttt{configs/dataset/}.

Before running experiments, set the data root so Hydra resolves paths correctly:

\begin{verbatim}
export DATA_ROOT=/path/to/project/data/processed
export OUTPUT_DIR=./outputs
export LOG_DIR=./logs
\end{verbatim}

\section{Training a Model}

Training is started from the repository root:

\begin{verbatim}
python -u src/train.py
\end{verbatim}

The default configuration is read from \texttt{configs/train.yaml}. It can be overridden from the command line. For example, to train a DeepLabV3+ model with a specific dataset and loss:

\begin{verbatim}
python -u src/train.py \
  dataset=sps_10_shadowed_2 \
  model=deeplabv3plus/resnet101_pretrained \
  loss=binarytanimoto \
  batch_size=16 \
  epochs=100
\end{verbatim}

Hydra creates a separate output directory for each run. The best checkpoint is saved as \texttt{best\_model.pth}. By default, the checkpoint is selected by positive-class \gls{iou}, but the configuration can use \gls{f1} instead through \texttt{checkpoint\_metric}.

\section{Evaluation}

Evaluation reconstructs a validation \gls{probabilitymap} from overlapping tiles, searches for an optimal threshold, and prints segmentation metrics. Run:

\begin{verbatim}
python -u src/evaluation.py \
  checkpoint_path=/path/to/best_model.pth \
  dataset=terraces_aniso_ladzany_128_32
\end{verbatim}

The evaluation output includes the optimal threshold and metrics such as positive \gls{iou}, mean \gls{iou}, Precision, Recall, \gls{f1}, specificity, accuracy, and \texttt{MOR10R}. The metrics are computed only over valid pixels and over regions actually covered by validation tiles.

\section{Full-Area Inference}

Inference generates a full \gls{probabilitymap} without requiring ground-truth labels. Run:

\begin{verbatim}
python -u src/inference.py \
  checkpoint_path=/path/to/best_model.pth \
  dataset=terraces_aniso_ladzany_full \
  output_path=outputs/inference/ladzany_probability_map.npy
\end{verbatim}

The output is a single-channel \texttt{.npy} array with probabilities in the range from 0 to 1. A \gls{binaryprediction} can be obtained by thresholding this array. The threshold may be the default value from \texttt{configs/inference.yaml}, a threshold found during evaluation, or a task-specific value chosen for visual \gls{gis} inspection.

\section{Running on a SLURM Cluster}

The repository contains \texttt{run\_slurm.slurm}, which demonstrates a cluster run. Before submitting it, update paths to the virtual environment and data directory. A typical submission is:

\begin{verbatim}
sbatch run_slurm.slurm
\end{verbatim}

The script exports \texttt{DATA\_ROOT}, \texttt{OUTPUT\_DIR}, and \texttt{LOG\_DIR}, activates the virtual environment, and starts \texttt{src/train.py} through \texttt{srun}.

\section{Correct Use and Limitations}

The pipeline assumes that all rasters, masks, and \gls{datasetmanifest} are spatially aligned. Incorrect alignment between the input raster and mask will produce invalid training data even if the scripts run without errors. For Slovak terraces experiments, the selected held-out location in the dataset configuration must match the intended leave-one-location-out fold.

The generated \gls{probabilitymap} should be interpreted as model outputs, not as verified archaeological discoveries. They should be checked against the original LiDAR-derived visualizations and, where possible, reviewed by a domain expert in a \gls{gis} environment.

\end{sloppypar}
